\subsection{Target benchmark and computational protocol}
\label{sec:data}

The final target benchmark in this study is a CL-20/HMX energetic-material dataset provided in DeepMD format and converted into a canonical extended-XYZ representation for training and evaluation. The current release contains two chemically distinct components:
\begin{itemize}
    \item a CL-20 subset with $4875$ labeled frames and $288$ atoms per frame; and
    \item an HMX subset with $9939$ labeled frames and $56$ atoms per frame.
\end{itemize}
Together they provide $14814$ force-labeled configurations. Each frame includes cell vectors, atomic coordinates, total energies, and per-atom forces. The data were staged within the collaboration from DeepMD-format labeled arrays into the canonical CL-20/HMX target extxyz file so that the same suite logic could be used across the public reference system, the public surrogate systems, and the final energetic-material benchmark.

The currently recoverable upstream package-level provenance is already more specific than an anonymous source note. The CL-20 subset maps to an upstream DeepMD root named \texttt{1-CL-20-Total} containing a single shard with $4875$ periodic $288$-atom configurations, whereas the HMX subset maps to \texttt{2-HMX-trainset3} containing two shards (\texttt{set.000} and \texttt{set.001}) with $5000$ and $4939$ periodic $56$-atom configurations, respectively.

Author-team confirmation from the BIT collaboration, including Dongping Chen as both a coauthor of the present manuscript and an author on the PCCP 2024 CL-20/HMX paper, now ties this packaged release directly to the CL-20/HMX data lineage reported by Wen \emph{et al.} \cite{wen2024pccp}. The staged HMX and CL-20 structures correspond to configurations at scale factors $0.96$, $1.00$, and $1.04$, with retained data at each scale for $300$, $500$, $1000$, $2000$, $3000$, and $4000$ K. Because the original DeepMD arrays were written in that fixed order, the retained subset can be mapped back to concrete temperature/scale conditions by frame order after conversion. The final subset was formed without relabeling. Instead, structures with any cell edge length below $1$ nm ($10$ \AA{}) were filtered out according to the collaboration team's benchmark-quality rule that periodic structures with such small cell edges should not enter the final benchmark. In practical terms, that filter removed roughly half of the original CL-20 payload and about one sixth of the HMX payload, with the HMX removals concentrated in the smallest-cell $300$ K subset.

The upstream AIMD and DFT labeling were performed in CP2K using the PBE generalized-gradient approximation, GTH pseudopotentials, DZVP-MOLOPT basis functions, and Grimme D3 dispersion correction, with auxiliary plane-wave cutoffs of $60$ Ry and $400$ Ry. That quality-control filter yields the present counts of $9939$ HMX and $4875$ CL-20 structures. In that sense, the underlying data payload is present within the author collaboration, the labels remain on the original published data lineage, and the package-to-extxyz conversion path is now locally verifiable. What remains incomplete is not the existence of the arrays themselves, but the reviewer-complete method record that would tie this packaged release to exact upstream shock-generation settings and a fully archived labeling workflow without ambiguity. The staged file preserves material identity and DeepMD shard identity in addition to the atomistic labels. In the current release, all frames carry the same source tag (\texttt{cl20\_hmx\_target\_v1}) and are labeled as \texttt{decomposition} configurations; the retained systems can be assigned to concrete temperature and scale conditions, but pressure and temperature metadata are not yet preserved explicitly as per-frame fields in the staged file.

The target split protocol is therefore conservative by construction. Because the current release does not preserve reusable multi-frame trajectory groups, the available \texttt{trajectory\_id} field is effectively frame-unique rather than trajectory-unique. We consequently use a deterministic contiguous frame-index split ($80\%/10\%/10\%$ for train/validation/test). An additional consequence of the present frame ordering is that the held-out test segment covers HMX only: all CL-20 frames lie within the training portion of the contiguous split. The present CL-20/HMX benchmark therefore supports static and nonequilibrium rollout comparisons on real energetic-material structures, but not a final pressure-, temperature-, or trajectory-history-resolved mechanistic decomposition or a trajectory-clean target split. At the same time, the static results now include a cleaner material-balanced split-v2 evaluation built from target-specific metadata, specifically to test whether the main CL-20/HMX static family ordering survives a jointly held-out CL-20/HMX mix rather than only the original HMX-only contiguous slice. To make these boundaries explicit, the submission package includes both a reviewer-facing benchmark-boundary table and a split-upgrade note that separate which target components already support the present claims from which metadata remain incomplete.

The target package should still expand further in future releases. In particular, an ideal production benchmark would also include:
\begin{itemize}
    \item trajectory-level pressure and temperature histories,
    \item reviewer-complete electronic-structure settings used for labeling, including exchange--correlation, basis or plane-wave convergence, and dispersion treatment \cite{perdew1996generalized,kresse1996efficient,grimme2010consistent},
    \item data-generation route for shock states, together with any reactive-force-field or accelerated sampling stages used to seed non-equilibrium configurations \cite{van2001reaxff},
    \item trajectory-level split rules that prevent leakage, and
    \item a regime-aware summary table for train, validation, and test subsets.
\end{itemize}

Within those current limits, the CL-20/HMX data already play a different role from the public surrogate systems. Transition1x and SPICE are used to characterize how charge-aware effects behave across publicly accessible chemical manifolds. The CL-20/HMX benchmark is the chemically decisive tier, because it asks whether the same modeling choices improve energetically relevant fitting and nonequilibrium stability on the actual energetic-material structures that motivate the paper, while also exposing which parts of that conclusion are presently bounded by missing provenance and thermodynamic metadata. In other words, the benchmark is already strong enough to test whether charge-aware modeling changes the static hierarchy and the frame-resolved failure structure on real energetic-material configurations, even though it is not yet strong enough to support a pressure-resolved or trajectory-history-resolved mechanistic attribution.
